\section{Discussion}
\label{sec:discussion}

\subsection{The DELB as an Integer-Valued Predictability Benchmark}
\label{subsec:discussion_rq1}

The first contribution of this study is an exact entropy-constrained MSE
lower bound for causal integer-valued prediction. Conditional Shannon
entropy is mapped through the inverse integer-lattice maximum-entropy
envelope, producing a model-independent benchmark for every admissible
integer-valued predictor under the specified information set and
squared-error loss.

The DELB should be interpreted carefully. It is exact in its discrete-envelope
form, but it does not
automatically equal the full Bayes risk. Universal attainment additionally
requires the prediction error to be independent of the admissible history
and to follow the moment-matched centered discrete-Gaussian distribution.
The predictor-dependent refinement accounts for residual dependence between
the error and the information set, while the two-slack decomposition
distinguishes this history-dependence component from error-distribution
shape mismatch.

More precisely, the comparison is with the infimum risk over measurable,
integer-valued, causal predictors having finite MSE under the stated
information set. The DELB does not exceed this optimal risk. If the
infimum is attained, equality holds only when an optimizing predictor
satisfies both attainment conditions. If no optimizer exists, numerical
equality of the two infima must not be described as attainment by a
predictor.

This distinction gives the DELB a diagnostic role beyond that of a scalar
benchmark. A large gap between achieved MSE and the universal DELB may
reflect remaining exploitable dependence, an unsuitable error distribution,
finite-sample model limitations, or some combination of these factors.
Consequently, the bound identifies a theoretically available region of
performance rather than guaranteeing that a practical estimator can reach
it.

The theorem is independent of the crime and bicycle application.
Numerical inversion and simulation verify the computational
implementation of the lattice envelope and its boundary cases, but the
theorem itself follows analytically. Likewise, the Fano analysis provides
a classification-loss counterpart and is not a second validation of the DELB.



\subsection{Finite-Sample Information Gain Requires Its Own Inferential Layer}
\label{subsec:discussion_rq2}

At the population level, adding side information reduces conditional
entropy by exactly its CMI with the target given the baseline information.
Monotonicity of the inverse lattice envelope then orders the corresponding
DELBs. This population statement is straightforward; estimating it from
sparse conditional states is not.

The known-truth simulations show that positive estimated CMI can
arise under zero population CMI. Adding a discretized variable divides each
baseline state into a larger number of joint states, many of which may be
weakly occupied or observed only once. Standard corrections reduce some
of this bias but do not remove the dependence of the estimate on state-space
geometry and sample size.

The estimator sensitivity analysis leads to the same caution. Plugin,
Miller--Madow, and observed-support Jeffreys--Dirichlet estimates produced
positive DELB reductions for all 16 prespecified specifications in each
city. Yet E17 found no rejection under any estimator, whether false
discovery rates were controlled within its nine city--null-design tests
or across all 27 estimator--city--design tests. The Jeffreys result is a
sensitivity check over observed support, not a reconstruction of
unobserved states.

A second distinction concerns randomization. A transformation-based
reference distribution is generated by repeatedly changing the observed
data structure and recomputing both the statistic and its support.
Its center therefore need not be zero and should not be interpreted as
either population CMI or a direct estimate of estimator bias. Valid
interpretation is specific to the estimator, transformation, sample,
and support produced by the randomization design.

The empirical support diagnostics reinforce this point. Effective
conditional degrees of freedom were strongly associated with the
reference-distribution center, and changes in support under randomization
were associated with differences between observed and reference CMI. These
findings do not provide a universal analytic bias formula, but they show
why comparing an estimated CMI directly with zero can be misleading
in sparse nested panels.


The resulting evidence hierarchy is therefore:
\begin{equation}
\begin{gathered}
\text{population DELB/CMI}
\longrightarrow
\text{finite-sample estimate}
\longrightarrow
\text{estimator behavior}\\
\longrightarrow
\text{design-conditioned reference}
\longrightarrow
\text{realized prediction error}.
\end{gathered}
\label{eq:discussion_inference_chain}
\end{equation}
Each stage answers a different question. A positive estimate describes
the observed partition; a randomization comparison evaluates a specified
null mechanism; and held-out prediction evaluates a particular fitted model.
None can substitute automatically for the others.


\subsection{What the Bicycle--Crime Analysis Does and Does Not Establish}
\label{subsec:discussion_rq3}

The urban results contain a strong descriptive pattern. Adding lagged
bicycle state reduced estimated conditional entropy and the corresponding
DELB in every city, every full-sample spatial--temporal scale, all 16
alternative theoretical specifications, and all exploratory crime categories.
The local estimates also showed spatial clustering in Washington, DC and
Vancouver. These results indicate that bicycle-augmented state partitions
consistently separate crime-count distributions more finely in the observed
samples.

The training-period-frozen analysis-domain comparison preserved the
direction of this pattern but changed its magnitude. Expanding from the
training-period bicycle-covered domain to the complete training-period
crime domain attenuated CMI and DELB reductions most strongly in New York
and Vancouver. The estimate therefore depends on the spatial target
population fixed before validation and testing; this domain contrast has
no causal interpretation.

The calibrated evidence was substantially weaker. None of the city-level
or local comparisons distinguished the observed CMI from all prespecified
temporal, spatial, and circular-shift references. Distant bicycle lags
yielded estimates comparable with the one-day lag, and none of the 270
multiscale randomization tests survived FDR correction. Thus, the present
design does not isolate an information contribution specific to recent,
geographically aligned bicycle mobility.

This finding should not be interpreted as proof that bicycle activity is
conditionally irrelevant to crime. The transformations are diagnostic
rather than exact conditional randomization procedures that preserve the
complete mobility distribution within every baseline state. Bicycle
use contains slow seasonal patterns, system-wide operational changes,
weather responses, pandemic-period shifts, and other common temporal
structure. These components can remain present in both recent and
transformed mobility series.

The baseline-state-matched restricted randomization reached the same
non-rejection while preserving grid and calendar strata more closely.
Its semi-synthetic Type I error was compatible with the nominal level, but
power remained low at CMI magnitudes comparable with the smaller empirical
signals. The added experiment therefore strengthens the calibration audit,
not a claim that mobility is conditionally irrelevant.

A more precise interpretation is therefore that the positive estimated
DELB reductions cannot, under the current calibration, be uniquely
attributed to recent bicycle-specific activity. This is a limit on
empirical attribution rather than a rejection of mobility information
in general.



\subsection{Theoretical Opportunity and Realized Prediction Are Different}
\label{subsec:discussion_prediction}

The held-out analyses illustrate why a population lower bound and a
fitted model's error must be evaluated separately. Adding information
can lower the best entropy-constrained error benchmark even when a
particular estimator cannot exploit that information from a finite sample.
Conversely, a model may improve by using seasonal or proxy structure
without establishing a bicycle-specific population CMI.

This distinction appeared clearly in the results. Bicycle-augmented DELBs
declined in all cities, whereas held-out improvements varied across models
and locations. Only two New York models reduced realized MSE by more than
the estimated DELB reduction and therefore narrowed the Predictability Gap.
The Fano comparison showed the same phenomenon under classification loss:
every estimated augmented floor declined, while the statewise modal
classifier worsened in all comparisons as unseen augmented states became
more frequent.

The Predictability Gap is consequently useful as a matched accounting
quantity rather than as an independent universal metric. Both MSE and
DELB must use the same target, forecast horizon, information set,
evaluation sample, and loss. A negative change in the gap does not
violate the lower bound; it indicates that the fitted model did not
realize the additional theoretical opportunity implied by the estimated
information state.

\subsection{Practical Implications for Urban Data and Forecasting Systems}
\label{subsec:practical_implications}

For urban analytics teams, the results support a staged data-value
assessment rather than a binary decision based on one fitted model.
Table~\ref{tab:practical_decision_framework} translates the evidence layers
into four operational questions. DELB first measures whether a candidate
source changes the estimated model-independent opportunity set. Support
diagnostics and randomization then assess whether that change is
distinguishable under explicit finite-sample references. Matched held-out
models finally determine whether an available forecasting system can use
the source in the deployment population.

\begin{table}[!htbp]
\caption{Practical decision framework for evaluating an added urban data
source in an integer-count forecasting system.}
\label{tab:practical_decision_framework}
\centering
\setlength{\tabcolsep}{3.5pt}
\begin{tabularx}{\textwidth}{
>{\RaggedRight\arraybackslash}p{2.75cm}
>{\RaggedRight\arraybackslash}X
>{\RaggedRight\arraybackslash}X
>{\RaggedRight\arraybackslash}p{3.15cm}}
\toprule
\textbf{Decision question} & \textbf{Required evidence} &
\textbf{Operational response} & \textbf{Boundary}\\
\midrule
Does the source create an estimated prediction opportunity? &
DELB and CMI point estimates in frozen target populations &
Retain the source for calibrated evaluation if the estimated floor
declines materially & Point estimates alone are insufficient for
source attribution or realized benefit\\
Is the apparent gain distinguishable from finite-sample structure? &
Support diagnostics, known-truth calibration, and design-matched
randomization & Attribute the gain only to the tested design and only when
calibration and power support separation & Low-powered non-rejection is
inconclusive; randomization provides no causal identification\\
Can available models use the source on unseen data? &
Matched held-out models with identical targets, samples, tuning rules, and
loss & Deploy only city--model combinations showing stable held-out benefit;
monitor the Predictability Gap & Model-specific gains require separate
validation across cities and model families\\
Does value persist outside the service footprint? &
Training-period-frozen coverage-conditioned and complete crime-supported
deployment populations &
Report both populations and treat coverage as part of the data product &
Service-area estimates apply only to covered locations\\
\bottomrule
\end{tabularx}
\end{table}


This sequence matters for procurement and engineering. A positive DELB
reduction can justify retaining a data source for further evaluation, but
not paying for, integrating, or operationalizing it on that basis alone.
In the present application, the large bicycle-domain reductions in New
York City and Vancouver attenuated in the complete crime-supported
populations, and the calibrated tests did not isolate a bicycle-specific
signal. Coverage, state support, and calibration should therefore be
treated as properties of the data product, not as post hoc sensitivity
checks.

Deployment decisions should also remain model and city specific. The
New York state-mean and gradient-boosting models narrowed the
Predictability Gap, whereas the remaining city--model combinations did
not. A practical pipeline should consequently preserve a baseline model,
evaluate the added source on a frozen holdout period, monitor both
\(\Delta\operatorname{MSE}\) and \(\Delta G\), and withdraw the added
feature when its benefit does not persist under rolling evaluation.
The DELB is useful here as an opportunity budget: it identifies how much
the estimated floor moved, while the gap shows whether the fitted system
converted that movement into realized accuracy.

These implications concern forecasting-system evaluation, not operational
policing decisions. The analysis is aggregate, noncausal, and based on a
mode-specific mobility proxy. It does not identify individuals, locations
that should receive enforcement, or the effect of bicycle activity on
crime. Any public-sector use would require separate fairness, privacy,
governance, and causal-impact assessment beyond the scope of the DELB
workflow.



\subsection{Methodological Contribution and Scope}
\label{subsec:discussion_contributions}

The theoretical contribution is not a new discovery of the discrete-Gaussian
extremizer, the lattice correction, the information-projection identity,
or Fano's inequality. It is the integration of established ingredients into
a four-part framework for causal integer-valued prediction:

\begin{enumerate}
\item an exact inverse-lattice MSE lower bound;
\item a predictor-dependent refinement;
\item necessary and sufficient attainment conditions; and
\item a two-component decomposition of nonattainment.
\end{enumerate}
The empirical contribution is the connection of this population theory
to a reproducible finite-sample evidence strategy. This connection is
important because entropy-based prediction limits are often estimated
precisely in settings where the state space is large relative to the
available sample. Without support diagnostics and design-matched calibration,
apparent reductions in a theoretical floor may be overstated as verified
information gain.

The DELB is narrower than general Bayes-risk or rate--distortion theory.
It assumes integer-valued targets, squared-error loss, and a discrete
representation of the conditioning state. This narrower scope permits
an explicit lattice envelope and interpretable equality conditions,
but it also makes entropy estimation vulnerable to sparse support.
The same trade-off is likely to arise in demand forecasting,
epidemiological counts, traffic incidents, service arrivals,
and other integer-valued systems.



\subsection{Limitations and Future Directions}
\label{subsec:limitations}

Several limitations concern the empirical data. Shared-bicycle trips are a
mode-specific trace of public-space activity rather than a census of
ambient population. Coverage depends on service boundaries, station
locations, adoption, system rebalancing, weather, endpoint completeness,
and recording practices. Crime data are also affected by victim reporting,
police classification, and jurisdiction-specific definitions. The 2020--2022
period includes pandemic-related structural change, and associations at
the 1 km grid--day level do not support individual-level inference.

Other limitations arise from the inferential design. The randomization
procedures do not preserve the exact conditional distribution of mobility
within every baseline state and therefore should be interpreted as
transformation-specific diagnostics rather than exact
conditional-independence tests. The known-truth simulation
spans a finite factorial family, and the support regressions
are descriptive diagnostics rather than analytic bias corrections.
The E18 semi-synthetic study used six representative grids, 365 days,
50 Monte Carlo datasets per cell, and 99 randomizations per dataset.
Its binomial intervals are correspondingly wide, several requested CMI
levels were not attainable on the fixed empirical support, and its low
power prevents interpreting empirical non-rejection as evidence of no
conditional information.
The pretest Fano distribution may differ from the test-period population.
For the multiscale randomization analysis below the full sample,
computation was restricted to one prespecified thinning path, although
estimator stability was examined over 50 paths.

Spatial and temporal aggregation also changes the target distribution.
Area--time normalization changes units but does not transform daily, weekly,
500 m, and 2 km outcomes into repeated measurements of one scale-free
population quantity. Claims of a universal scale law would require a
separate theoretical definition and study design.

Future work could reduce state sparsity through shrinkage, hierarchical
pooling, Bayesian or frequentist partial pooling, and cross-fitting.
More conditional surrogate designs could preserve mobility distributions
within baseline states more closely. Empirical extensions should
incorporate additional mobility modes, pedestrian activity, transit use,
weather, and post-2022 observations while preserving matched causal
information sets. Theoretical extensions include sharper multivariate
lattice envelopes, alternative integer-valued loss functions, and bounds
for dependent vector count processes.





\section{Conclusions}
\label{sec:conclusions}

This study formulated the Discrete Entropy Lower Bound for causal
integer-valued prediction. By inverting the exact integer-lattice
maximum-entropy envelope, the DELB maps conditional Shannon entropy
to a model-independent, entropy-constrained MSE lower bound.
The framework further provides a predictor-dependent refinement,
necessary and sufficient attainment conditions, and a two-component
decomposition that separates residual history dependence from
error-distribution shape mismatch.

The study also established a finite-sample evidence strategy for
evaluating changes in this bound after side information is added.
Population CMI, estimated CMI, estimator bias, transformation-specific
randomization references, and achieved prediction error are distinct
quantities. Known-truth simulations and support diagnostics showed that
sparse nested states can produce positive CMI estimates and nonzero
randomization centers even under a population null, making design-matched
calibration essential.

In the three urban panels, lagged bicycle information consistently
lowered estimated DELBs across cities, spatial and temporal scales,
crime categories, and alternative specifications. However,
the observed reductions were not distinguishable from the prespecified
temporal, spatial, or circular-shift references after calibration, and
they did not translate uniformly into lower held-out prediction error.
The findings therefore do not establish a recent bicycle-specific
reduction in the population prediction bound, nor do they imply that
bicycle mobility is predictively irrelevant. Instead, they illustrate
the central methodological distinction of this study: a lower estimated
uncertainty bound, calibrated evidence of added information, and improved
realized prediction are separate claims requiring separate evidence.
More broadly, the DELB provides a rigorous framework for studying
information-theoretic prediction floors and finite-sample information
value in crime counts and other integer-valued stochastic systems. In
practice, it supports a staged decision process in which estimated
opportunity, calibrated separation, and realized held-out benefit must be
established separately before an added urban data source is deployed.
